\section{Methods description form (Track 2 methods)}
\textit{Answers to the organisers' official methods description template (version with the required LLM question), reproduced here because the submission form has no separate upload field for it.}

\begin{description}[leftmargin=0pt,style=nextline,itemsep=4pt,font=\sffamily\bfseries\small\color{branddark}]
\item[Team name] aperezcastro (Agustín Pérez Castro, individual)
\item[Please describe your approach in detail (how you went from variant/mechanism to candidate medication(s)).] Allele-aware, patient-aware pipeline (scripts/t2/). (1) Mechanism model from Track 1: null allele p.(Leu737*) (UGA, NMD) + kinase-domain missense p.(Asn1002Lys) $\rightarrow$ partial loss of function in a dosage-sensitive checkpoint. The CLASS of the missense allele (low abundance vs stable-but-dysfunctional) is unmeasured; every disease-associated substitution tested in or near the kinase domain was 5-10x less abundant and the known stable alleles lie outside that domain, so we carry both scenarios at a 0.75/0.25 prior and price the uncertainty into mechanistic fit, not evidence. AlphaFold structure: Asn1002 buried high-confidence residue near the MVA1 allele L1012P (destabilisation). (2) Knowledge base (hipotesis.yaml): 7 therapeutic lines + a safety card, each with PMID-cited evidence graded by level and direction (supporting/contrary): H1 stabilise BUBR1 (NAD+/SIRT2), H2 read-through/NMD inhibition of the nonsense allele, H3 mTORC1 (sarcopenia), H4 senolytics, H5 aneuploid-selective stress, H6 interferon/JAK-STAT (biomarker-gated), H7 proteostasis amplification/antioxidants (evaluated and declined). (3) Connectivity (cmap.py): consensus of 32 BUB1B CRISPR-knockout L1000 signatures (20 cell lines) via SigCom LINCS, two-sided rank enrichment against chemical perturbation signatures, cross-cell-line consistency, Enrichr interpretation (interferon/NF-kB/p53 up, E2F/G2M down). Used to check H6 for class consistency (JAK inhibitors reverse the signature), to support the safety card (DNA-damaging drugs mimic it) and to challenge H3 (mTOR inhibitors mimic it). (4) candidatos.py: seeded approved drugs (ChEMBL target expansion is available but added none in this run), annotated with ChEMBL and openFDA labels (paediatric use, boxed warnings) and penalised by the patient's HPO-derived risk profile; transparent score 0.30 evidence + 0.25 mechanistic fit + 0.20 paediatric feasibility + 0.25 patient-specific safety, with Dirichlet weight sensitivity; only approved medicines are ranked (discussion \#5). (5) Clinical tiering by benefit-risk: tier 0 safety note plus rung-0 measurement; tier 1 intermittent dasatinib (highest-scoring approved drug), to be tested first in the BubR1 H/L1002P mouse; tier 2 an NAD+ precursor - niacin (nicotinamide ties on score but inhibits sirtuins) - gated on the allele class AND documented remission; tier 3 low-dose sirolimus for muscle with \%MN-RET as a safety readout; tier 4 JAK inhibitors only if an interferon signature is measured; read-through closed on current allele-level evidence; proteostasis declined because raising HSF1 is a documented aneuploidy-TOLERATING event (Donnelly 2014). (6) nof1.py: an n-of-1 power simulation on micronucleated reticulocytes (\%MN-RET), the only endpoint for this disease that moves in weeks. Validation ladder: rung 0 measure \%MN-RET and its variability; patient cells (decides the allele scenario); iPSC myotubes; the BubR1 allelic mouse; then an n-of-1 protocol with pre-declared responder criteria.
\item[(required) If commercially-available Generative AI was used, record the provider, the plan or tier, and the relevant setting in your methods description. A line is enough, for example, "Anthropic API, Claude Sonnet, commercial terms, no training on customer content."] Anthropic Claude via the Claude Code CLI (Claude Opus models), Claude Max (consumer tier), model-training opt-out enabled ('Help improve Claude' off), no feedback provided on any output, no content flagged for safety review - per the organisers' confirmed guidance of 28 Sep 2026 (discussion \#2). Applicable terms: opted-out conversations are not used for training and are retained 30 days. Additional data minimisation: the model only saw code, aggregate counts, HPO terms and short candidate lists ($\leq$ 50 annotated variants); raw data and genome-scale derivatives were blocked from the agent by configuration. No patient variant was sent to any third-party API.
\item[Was candidate identification automated (e.g., via a drug-target database query or LLM-assisted search), or primarily manual literature review and expert curation?] Hybrid: candidate expansion, annotation, safety flags, connectivity and scoring are automated; the hypotheses, evidence grading, mechanistic-fit values and final clinical tiering were curated manually (LLM-assisted literature search with every reference verified against PubMed).
\item[Please describe any manual literature review or expert curation in detail.] PubMed searches at gene/pathway/drug level only (no patient data). Each claim in hipotesis.yaml cites a PMID whose title/abstract was retrieved and checked; contrary evidence is recorded explicitly (e.g., SIRT2 does not extend lifespan in wild-type mice; aneuploid-cell immune clearance argues against JAK inhibition). Mechanistic-fit values (0-1) are justified in the YAML. Automated label flags were reviewed manually.
\item[Please state whether you relied solely on publicly available drug/target databases and literature (preferred), or whether proprietary data sources were also used.] Solely public databases and literature.
\item[Please describe the public data sources used in detail (e.g., DrugBank, ChEMBL, ClinicalTrials.gov, OMIM, PubMed literature search).] PubMed/NCBI E-utilities; ChEMBL REST API (molecules, mechanisms, targets); openFDA drug labels; ClinicalTrials.gov API v2; LINCS L1000 via SigCom LINCS (l1000\_xpr knockout and l1000\_cp chemical signatures); Enrichr (MSigDB Hallmark 2020); UniProt O60566; AlphaFold DB AF-O60566-F1 v6; OMIM/HPO.
\item[Please describe any proprietary data sources used in detail.] None.
\item[How did you characterize the variant's mechanism (e.g., loss-of-function vs. gain-of-function, pathway disrupted, downstream biological consequence)?] Loss of function by reduced dosage: the nonsense allele is predicted to undergo NMD (no protein). For the kinase-domain missense we DO NOT assert a class: Suijkerbuijk 2010 splits MVA alleles into low-abundance and functionally-impaired classes, and although all disease-associated substitutions tested in or near the kinase domain were 5-10x less abundant, this residue has never been characterised, so scenario A (0.75) and scenario B (0.25) are both carried and a BUBR1 immunoblot in patient cells is the pre-declared gate. Pathway: spindle assembly checkpoint / mitotic checkpoint complex (APC/C-CDC20 inhibition) and kinetochore-microtubule attachment. Downstream: chromosome missegregation $\rightarrow$ micronuclei $\rightarrow$ cGAS-STING interferon/NF-kB and p53 responses (confirmed in the knockout signature), senescence/SASP, mTORC1 hyperactivity in muscle (mouse allelic series), tumour predisposition.
\item[Please provide an estimate of time or effort spent on this analysis (if possible).] About 2 working days of analysis and writing; computation runs in minutes on a laptop (public APIs, cached).
\item[Provide a method abstract (up to 500 words). Please include strengths and limitations of your approach as well as methodological detail.] We present an allele-aware, patient-aware repurposing pipeline for a child with biallelic BUB1B mosaic variegated aneuploidy (null p.Leu737* + kinase-domain p.Asn1002Lys). The mechanism is partial loss of function in a dosage-sensitive checkpoint; the missense allele class is unmeasured and both scenarios are carried. Three evidence streams were combined: (i) a PMID-cited knowledge base of seven therapeutic lines plus a safety card, scoring supporting and contrary evidence explicitly and weighting target validation below demonstrated benefit; (ii) connectivity analysis of 32 BUB1B-knockout L1000 signatures across 20 cell lines; (iii) automated annotation from ChEMBL and FDA labels, penalised by the patient's HPO-derived risk profile, with Dirichlet weight-sensitivity analysis. The knockout signature is interferon/NF-kB/p53-up and E2F/G2M-down; DNA-damaging chemotherapeutics reproduce it and mTOR inhibitors mimic it; JAK inhibitors reverse it with class consistency but at ranks 15-154 of 539 reversers, a targeted search rather than a discovery. Two evidence reviews changed most lines: read-through was closed, on current evidence, by the stop-codon context (UGA, unfavourable +4 A, NMD) plus evidence that ataluren's activity is a luciferase artefact; nicotinamide fell because rhabdomyosarcoma depends on NAMPT/NAD+ and nicotinamide inhibits sirtuins; senolysis became a preclinical experiment after the best senolytic RCT missed its endpoint; and sirolimus fell from first to tenth when its mouse support proved to come from monoallelic carriers. THE MAIN RESULT is a decision rule: classifying every line by what it does to a chromosomally unstable cell (correct fewer errors / clear the abnormal cells / tolerate help them cope) shows that no axis wins on evidence and that each corrective line is blocked by a different nameable fact about this child. Today the best-supported option is intermittent dasatinib (clear axis, narrowly ahead of the JAK inhibitor ruxolitinib); the mechanism-matched candidate, an NAD+ precursor (niacin), ranks eighth, and rises to first if the clinical team establishes the allele class and documented remission - so the first recommendation is a measurement, not a prescription. In a cancer-predisposition syndrome a premalignant clone is an abnormal cell too. Deliverables: a two-page pharmacogenomic safety card, usable now; \%MN-RET, a blood measurement never made in MVA, diagnostic first and then an endpoint with its power calculation; a ranked and conditioned candidate list, and a four-rung validation ladder (patient cells, iPSC myotubes, the published BubR1 allelic mouse model, then an n-of-1 protocol with pre-declared responder criteria) built on verified precedents (lonafarnib, sirolimus in Castleman disease). Strengths: per-allele strategies, explicit contrary evidence that demoted our own candidates, patient-specific safety beyond drug labels, and reproducible gene-agnostic code. Limitations: efficacy evidence is preclinical; connectivity signatures come from tumour cell lines after acute knockout; mechanistic-fit values are expert judgements and the ranking is sensitive to two of them, so the counterfactual is published rather than a robustness claim; every ranked drug was hand-seeded (automated target expansion added none); the \%MN-RET baseline and variability used in the power calculation have never been measured in MVA; regulatory routes are orientation, not legal advice. An independent blind adversarial review of an earlier draft drove these corrections.
\end{description}
